\chapter*{Abstract ...short}

\addcontentsline{toc}{chapter}{Abstract}

This work introduces Alignment-Based Information Modelling (AIM),
a unified framework for railway alignment representation, analysis,
and optimization.

Alignment is formulated as a state-based object combining geometry,
curvature, and cant, enabling direct evaluation of dynamic quantities
and supporting numerical reconstruction and constrained optimization.

Transition curves are reinterpreted as curvature control functions and
derived as solutions of variational problems.
A normalization pipeline allows integration of heterogeneous data
sources such as IFC and LandXML.

The concept of BIMalignment is proposed as an extension of BIM,
establishing alignment as the central organizing structure of
infrastructure models.

The framework provides a foundation for dynamic, optimization-driven,
and computationally integrated railway alignment design.
